\href{https://www.archives.gov/milestone-documents/executive-order-9066}{Executive Order 9066} 
was signed on February 19, 1942 and granted the Secretary of War the authority
to designate military areas from which civilians could be prevented access. This
order was used to target all persons of Japanese descent living in the so-called
`Evacuation Zone' which included all of California, and large areas of
Washington, Oregon, and Arizona shown in blue in Figure \ref{fig:internmentmap}.
At first, the Secretary of War Henry Stimson encouraged some families to relocate voluntarily from the West Coast,
but starting at the end of March 1942, 
the Wartime Civil Control Administration (WCCA) and War Relocation Authority (WRA)
were established to remove all remaining Japanese Americans to temporary
Assembly Centers which were makeshift shelters in close proximity to the
population centers where these people were living. As the more permanent WRA
Relocation Centers reached completion in May 1942, all Japanese Americans left
in the West Coast were moved from Assembly Centers to Relocation Centers by
train. This process lasted until November of 1942
\citep{commissiononwartimerelocationPersonalJusticeDenied1983}.

After two years of WRA camp operations, the government began to debate how to
end the evacuation orders towards the end of 1944. One statement issued on
December 9 illustrates that one of the initial goals of interment was actually
to cause a permanent shift in the Japanese American population away from the
West Coast as part of a wide-scale program of `deghettoization':
\begin{quote}
  ``One of the major WRA aims, from the beginning, has been to encourage the widest possible dispersal of evacuees throughout the Nation, and this will continue as a prime objective during the final phase of the program.''
  \hfill \cite[p.~234]{commissiononwartimerelocationPersonalJusticeDenied1983}
\end{quote}

As seen in Figure \ref{fig:1940JAmap}, over 90\% of the 1940 Japanese American
population lived in areas which would become subject to relocation orders.
This population was overwhelmingly concentrated along the West Coast, particularly in California,
which was home to large, well-established Japanese American communities,
many of whom were engaged in farming and small businesses.
The geographic concentration of this population in economically vital areas near major cities and military installations heightened the perceived security threat,
prompting their forced relocation.
